\chapter*{第一卷说明}
\addcontentsline{toc}{chapter}{第一卷说明}
本卷依次包括第 I-01 章《预备知识与依赖边界》、第 I-02 章《度量空间、完备性与压缩原理》、第 I-03 章《赋范空间与 Banach 空间》、第 I-04 章《Hilbert 空间的几何》、第 I-05 章《有界线性算子与连续线性泛函》、第 I-06 章《Hahn--Banach 与分离》、第 I-07 章《对偶、二次对偶与弱/弱*拓扑基础》、第 I-08 章《Baire 纲定理与一致有界性》、第 I-09 章《开映射、有界逆与闭图像》、第 I-10 章《弱紧性、Banach--Alaoglu 与自反性》、第 I-11 章《对偶算子、Hilbert 伴随与算子拓扑》、第 I-12 章《预解集、谱与最小 Banach 代数工具》、第 I-13 章《紧算子、Riesz--Schauder 与 Fredholm 择一原理》和第 I-14 章《紧自伴谱定理与第一遍闭合》. 内容从拓扑与度量基础出发，经 Banach 与 Hilbert 空间、有界线性算子、Hahn--Banach 理论、弱拓扑与 Baire 纲方法，建立一致有界性、开映射、有界逆、闭图像、Banach--Alaoglu、Goldstine、自反空间弱紧性、对偶算子、Hilbert 伴随、强/弱算子拓扑，以及预解集、谱、谱半径、多项式谱映射、紧算子、Riesz--Schauder 结构、Fredholm 择一原理与紧自伴谱定理等第一卷线性主干，并以正交特征分解与离散函数演算完成第一遍闭合.

本卷固定 \(\displaystyle \mathbb N=\{1,2,\dots\}\)、\(\displaystyle \mathbb N_0=\{0,1,2,\dots\}\) 与 \(\displaystyle \mathbb K\in\{\mathbb R,\mathbb C\}\). 弱紧性与弱*紧性均按一般拓扑紧性理解，并以网/子网作为基本刻画；Eberlein--\v Smulian 定理尚不进入当前证明链. I-11 已严格分流 Banach 对偶算子与 Hilbert 伴随并建立强/弱算子拓扑；I-12 已在非零复 Banach 空间上建立预解集、谱非空紧性与谱半径公式，并以显式复化处理实空间接口；I-13 已建立紧算子、Riesz--Schauder 理论与 Fredholm 择一原理；I-14 已在实、复 Hilbert 空间上建立紧自伴谱定理、正交特征分解与离散函数演算. 一般有界正规谱定理留到第二卷，弱下半连续性、强制性与变分直接法留到第三卷相应章节. 对任意弱收敛网，不断言其整个取值范围范数有界.
